% Complete source sections 12--13 plus their required division, source-algebra
% and Sobolev-comparison proofs. Source SHA256:
% b70272d1e23d49d2c5b7c9d407a01b151443a17ea93a7e6a3f4597e0e982f90a.
% Integration changes only namespace, exact local cross-references and the
% explicit time-first field convention; dependency additions are recorded.
\section{The actual source quotient of a flow-driven arithmetic profile}
\label{csq:sec-flow-source-quotient}


\subsection{Retained source data and the full division estimate}
The arithmetic data are exactly those of
Section~\ref{sec:cq-source} and \eqref{eq:cq-N}:
$\cqE$ is the set of entire functions of order at most one,
$\cqN=\overline{\mathbb C[s]\Xi}^{\,\tau}$,
$\cqQ=(\cqE,\tau)/\cqN$, and $S[h]=[sh]$.
The source topology $\tau$ remains unnamed beyond the primary
specification; it is not replaced by the explicit growth topology
$\gamma$ of Section~\ref{sec:cq-heatspace}.
Its growth seminorms are
$\|h\|_{p,a}=\sup_s|h(s)|e^{-a|s|^p}$ for every $1<p<2$, $a>0$.
The full-domain heat group is precisely \eqref{eq:cq-U}.
The original multiplication $M h(s)=s h(s)$ preserves $\cqN$,
as proved from its polynomial generators in Section~\ref{sec:cq-source}.
In the flow formulas below, $S_j$, $J$ and $X_t$ instead denote
the polynomial target, its derivative and the volume-preserving
preterminal material flow of Section~\ref{sec:ns-full-flow};
their exact formulas are \eqref{nfl:S}, \eqref{nfl:Jinverse} and
\eqref{nfl:compact}. The field convention is $u(t,x)$.
The particular source field and its full remainder are those of
Section~\ref{sec:ns-actual-witness}. Its cited existence theorem
remains an imported primary-source result.

For the division argument the value $\Xi(0)=\xi(1/2)$ is nonzero.
One direct verification retains the usual Dirichlet eta function:
$\eta(1/2)=\sum_{n\geq1}(-1)^{n-1}n^{-1/2}>0$, since its paired
terms are positive and its first pair is strictly positive.
The alternating-series convergence and
$\eta(s)=(1-2^{1-s})\zeta(s)$ give
$\zeta(1/2)=\eta(1/2)/(1-\sqrt2)<0$.
The other factors in
$\xi(1/2)=\tfrac12(1/2)(-1/2)\pi^{-1/4}\Gamma(1/4)\zeta(1/2)$
therefore give $\Xi(0)>0$. No assertion about any noncentral
zeta zero is used.

We now compute the closures and the failure of descent in specified
topologies. This does not identify them with $\tau$.
Let $\kappa$ be the topology of uniform convergence on compact sets
on the same set $\cqE$, and define
\[
 N_\gamma=\overline{\cqC[s]\Xi}^{\,\gamma},\qquad
 N_\kappa=\overline{\cqC[s]\Xi}^{\,\kappa}.
\]
Evaluations of all derivatives are continuous in both topologies by
Cauchy's formula. The growth topology is finer than $\kappa$, so
$N_\gamma\subseteq N_\kappa$.
If $\lambda$ is a zero of $\Xi$ with its actual multiplicity $m_\lambda$,
every element of either closed subspace has vanishing derivatives
of orders $0,\ldots,m_\lambda-1$ at $\lambda$.

\begin{cqlemma}[Entire division with its growth and topology retained]
\label{csq:thm-division}
Let $g\in\cqE$ with $g(0)\ne0$. If $f\in\cqE$ and $h=f/g$
extends to an entire function, then $h\in\cqE$. For $1<p<2$,
$a>0$, put $b=a/(6\cdot2^p)$. There is the linear division estimate
\begin{equation}\label{csq:eq-division-bound}
 \|f/g\|_{p,a}\leq
 \left(\frac{\max(1,\|g\|_{p,b})}{|g(0)|}\right)^3
 \|f\|_{p,b}.
\end{equation}
Multiplication by $g$ is a topological linear isomorphism from
$(\cqE,\gamma)$ onto the closed subspace of functions divisible by
$g$, with its subspace topology. Compact-topology division is also
continuous.
\end{cqlemma}
\begin{proof}
Choose $R>0$ with no zero of $g$ on $|s|=R$. Factor its finitely
many zeros in the disk, retaining multiplicities. The residual is
holomorphic and zero-free there, so its log modulus is harmonic.
For $|z|<R$ the boundary mean of $\log|Re^{i\theta}-z|$ equals
$\log R$: expand $\log(1-(z/R)e^{-i\theta})$ in its uniformly
convergent series and take real parts. The mean-value formula for
the residual gives the exact Jensen identity
\[
 \frac1{2\pi}\int_0^{2\pi}\log|g(Re^{i\theta})|\,d\theta
 =\log|g(0)|+\sum_{|z|<R}m_z\log(R/|z|)\geq\log|g(0)|.
\]
There is no zero at zero. With
$\log^+x=\max(0,\log x)$ and $\log^-x=\max(0,-\log x)$,
the inequality $\log^+|h|\leq\log^+|f|+\log^-|g|$ gives
\begin{equation}\label{csq:eq-division-mean}
 \frac1{2\pi}\int_0^{2\pi}\log^+|h(Re^{i\theta})|\,d\theta
 \leq\log^+M_f(R)+\log^+M_g(R)-\log|g(0)|.
\end{equation}
The right side is nonnegative since $M_g(R)\geq|g(0)|$.
The harmonic Poisson extension of the continuous nonnegative boundary
function $\log^+|h|$ majorizes $\log|h|$ by the maximum principle
and is nonnegative; hence it majorizes $\log^+|h|$ in the disk.
For $|s|\leq r<R$ the Poisson kernel is at most $(R+r)/(R-r)$.
Take $R\downarrow2r$ through circles avoiding zeros. Continuity of
the two maximum functions gives, for every $r>0$,
\begin{equation}\label{csq:eq-division-radial}
 \log^+M_h(r)\leq
 3\bigl(\log^+M_f(2r)+\log^+M_g(2r)-\log|g(0)|\bigr).
\end{equation}
No pointwise lower bound for $g$ on a large circle has been assumed.
For $f\ne0$, put $F=\|f\|_{p,b}>0$ and apply this to $f/F,h/F$.
Then $\log^+M_{f/F}(2r)\leq b(2r)^p$, and
$\log^+M_g(2r)\leq\log\max(1,\|g\|_{p,b})+b(2r)^p$.
The coefficient of $r^p$ is $6b2^p=a$, giving
\eqref{csq:eq-division-bound}; continuity fills in $r=0$.
The zero numerator is immediate. All $p,a$ are allowed, so $h$
belongs to the original order class and division is continuous.
Multiplication is continuous by Lemma~\ref{thm:cq-density}.
The requisite vanishing jets at every zero are closed conditions;
local Taylor division identifies them with entire divisibility.
The estimate proves this closed set is exactly $g\cqE$.
For the compact topology choose $R>r$ with no boundary zero and
$\delta=\min_{|s|=R}|g(s)|>0$. The maximum principle for the filled
entire quotient gives $M_h(r)\leq\delta^{-1}M_f(R)$, proving
compact continuity too.
\end{proof}

For general nonzero $g\in\cqE$, choose and retain a point $c$ with
$g(c)\ne0$ and use disks centered at $c$. The functions and spectral
coordinate remain unchanged. The bound
$|s+c|^p\leq2^{p-1}(|s|^p+|c|^p)$ proves continuity of translations
by $c$ and $-c$ in $\gamma$; compact containment proves it in
$\kappa$. Thus all conclusions hold with the stated translated
estimates. The original center $c=0$ already applies to $\Xi$.

\begin{cqproposition}[The computed closures are the same subspace]
\label{csq:thm-equalclosures}
With their definitions unchanged,
\begin{equation}\label{csq:eq-equalclosures}
 \begin{aligned}
 N_\gamma=N_\kappa&=\Xi\cqE\\
 &=\{f\in\cqE:D^hf(\lambda)=0\text{ for every }\Xi(\lambda)=0,
                                      \ 0\leq h<m_\lambda\}.
 \end{aligned}
\end{equation}
This is equality of subsets of the original ring, not an
identification of $\gamma$, $\kappa$, or $\tau$.
\end{cqproposition}
\begin{proof}
Jet continuity puts both closures inside the displayed vanishing
set and gives $N_\gamma\subseteq N_\kappa$. Local Taylor division
makes $h=f/\Xi$ entire for any $f$ in that set; the division lemma
then puts $h$ in $\cqE$. Its Taylor polynomials $p_n$ converge in
$\gamma$, and multiplication gives $p_n\Xi\to f$ there. Thus the
vanishing set is contained in $N_\gamma$, proving all equalities.
The algebraic-ideal equality is the conclusion of this argument;
it was not substituted for an unevaluated closure.
\end{proof}


\subsection{Translation, heat, and the original spectral generator}

For $a,\theta\in\mathbb C$ define, with the input coordinate $s$ and
output coordinate $\zeta$ retained,
\[
 (T_ah)(\zeta)=h(\zeta+a),\qquad
 V_{\theta,a}=T_aU_\theta:\cqE\longrightarrow\cqE.
\]
The derivative is the complex derivative in the displayed coordinate.
Translation preserves the original order class. More precisely, for
$1<p<2$ and $c>0$ the already defined growth seminorms satisfy
\[
 \|T_ah\|_{p,c}\leq
 e^{c|a|^p}\|h\|_{p,c\,2^{1-p}},
\]
because $|\zeta+a|^p\leq2^{p-1}(|\zeta|^p+|a|^p)$.
Thus $T_a$ and $T_{-a}$ are continuous for $\gamma$. Ordinary
differentiation gives $DT_a=T_aD$; continuity and the convergent
heat series give $T_aU_\theta=U_\theta T_a$. Consequently
\[
 V_{\theta,a}^{-1}=U_{-\theta}T_{-a},\qquad
 V_{\theta,a}D=DV_{\theta,a}.
\]
These assertions hold on all of $\cqE$ and do not impose continuity
of translation or differentiation for the unnamed source topology.

Retain the actual $\tau$, $\cqN$ and $\cqQ$, and define their exact
transported presentations
\begin{equation}
 \tau_{\theta,a}=(V_{\theta,a})_*\tau,\qquad
 N_{\theta,a}=V_{\theta,a}\cqN,\qquad
 Q_{\theta,a}=(\cqE,\tau_{\theta,a})/N_{\theta,a}.
 \label{csq:eq-flow-moving-source}
\end{equation}
Then $N_{\theta,a}$ is closed, and
\[
 \overline V_{\theta,a}:\cqQ\longrightarrow Q_{\theta,a},
 \qquad[h]\longmapsto[V_{\theta,a}h]
\]
is a complex-linear homeomorphism. Its inverse is induced by the
displayed inverse on representatives. The original arithmetic
multiplication has the exact conjugate
\begin{equation}
 B_{\theta,a}=V_{\theta,a}sV_{\theta,a}^{-1}
              =\zeta+a-\frac{\theta}{2}D_\zeta .
 \label{csq:eq-flow-spectral-generator}
\end{equation}
Indeed the previous heat commutator gives
$U_\theta sU_{-\theta}=s-\theta D/2$; translation sends multiplication
by $s$ to multiplication by $\zeta+a$ and commutes with $D$.
This full combination preserves $N_{\theta,a}$ and is continuous for
$\tau_{\theta,a}$ by its conjugacy. The induced quotient operator is
similar through $\overline V_{\theta,a}$ to the source $S$, so its
spectrum is exactly the original $s$-spectrum. To verify both
directions, conjugate a continuous two-sided inverse of $S-\lambda$
by $\overline V_{\theta,a}$, or conjugate an inverse of the transported
operator back. At $\theta=0$, multiplication by $\zeta$ itself
descends, and its spectrum is $\operatorname{Spec}(S)-a$ because
$B_{0,a}=\zeta+a$.

Along the retained real flow let
\[
 a_j(t,q)=S_j(X_t(q)),\qquad
 v_j(t,q)=\partial_ta_j(t,q),\qquad
 \Psi_j(t,q,\zeta)=V_{\vartheta(t),a_j(t,q)}\Xi(\zeta).
\]
Here $S_j$ denotes the original polynomial target component, whereas
the unindexed $S$ above denotes the arithmetic quotient operator.
The source quotient is taken fibre by fibre at each fixed $(t,q,j)$.
If a topology on the whole family of quotient fibres is needed, give
their disjoint union the topology transported from
$\mathbb C^2\times\cqQ$ by
$((\theta,a),[h])\mapsto((\theta,a),[V_{\theta,a}h])$.
This specifies joint continuity by an exact trivialization and does
not assert joint parameter continuity for the untransported $\tau$.

Because $\Xi\in\cqN$, every profile represents zero in its actual
transported quotient:
\begin{equation}
 [\Psi_j]_{Q_{\vartheta,a_j}}=0.
 \label{csq:eq-flow-profile-zero}
\end{equation}
Its temporal defect is nevertheless a specific full representative,
\[
 \mathcal D_j=\partial_t\Psi_j+
              \frac{\vartheta'}4D_\zeta^2\Psi_j
             =v_jV_{\vartheta,a_j}\Xi'.
\]
This is the complete chain rule with the original heat coefficient.
Write
\[
 \beta=[\Xi']_{\cqQ},\qquad
 \mathfrak a_\beta=\{c\in\mathbb C:c\Xi'\in\cqN\}.
\]
Then the exact quotient statement and its scalar kernel are
\begin{equation}
 [\mathcal D_j]_{Q_{\vartheta,a_j}}
    =\overline V_{\vartheta,a_j}(v_j\beta),\qquad
 \ker(c\mapsto c\beta)=\mathfrak a_\beta .
 \label{csq:eq-flow-source-beta}
\end{equation}
The set $\mathfrak a_\beta$ is a closed complex linear subspace:
it is the inverse image of the closed $\cqN$ under the continuous
scalar map $c\mapsto c\Xi'$ into $(\cqE,\tau)$.
Its equality with $\{0\}$ or $\mathbb C$ is precisely the unresolved
membership question $\Xi'\notin\cqN$ or $\Xi'\in\cqN$.
No evaluation functional on the source quotient is used to decide it.

At a fixed physical point $x=X_t(q)$, retain
$v=J(x)u(t,x)$ and the full inverse
$J^{-1}(x)=[2W_1(x)\ W_2(x)\ W_3(x)]$. Thus the exact kernel of
the complexified velocity-to-source-class map is
\[
 J(x)^{-1}(\mathfrak a_\beta^3)\subseteq\mathbb C^3;
\]
for real velocities take its intersection with $\mathbb R^3$.
This follows componentwise from \eqref{csq:eq-flow-source-beta} and
the invertibility of the unchanged $J$.

The derivative relation explains why a zero profile class does not
determine its representative derivative. At a fixed $(\theta,a)$ put
\[
 \mathfrak D_{\theta,a}
 =\{([h],[Dh]):h\in\cqE\}\subset Q_{\theta,a}^2,\qquad
 \mathcal W_{\theta,a}
 =(DN_{\theta,a}+N_{\theta,a})/N_{\theta,a}.
\]
It has full domain and output fibre $[Dh]+\mathcal W_{\theta,a}$
at $[h]$: replace $h$ by every $h+n$ with $n\in N_{\theta,a}$.
Its presenting map has kernel
$\{h\in N_{\theta,a}:Dh\in N_{\theta,a}\}$.
Since $D$ commutes with $V_{\theta,a}$, conjugating both components
by $\overline V_{\theta,a}$ gives exactly the previous source
derivative relation and its ambiguity space. In particular
$\overline V_{\theta,a}\beta$ is an actual element of its
multivalued part at zero. The displayed temporal connection is an
identity on the full smooth profiles; a derivative on all
untransported source quotient representatives is not presumed.

\subsection{An exact information-preserving refinement and its residue projection}

Retain the explicit analytic ideal $I=\Xi\cqE$, already proved closed
for $\gamma$ by the full division estimate, and put
\begin{equation}
 H_\cap=\cqN\cap I,\qquad
 C_\cap=(\cqE,\tau\vee\gamma)/H_\cap,\qquad
 Q_I=(\cqE,\gamma)/I .
 \label{csq:eq-flow-meet-quotient}
\end{equation}
The join topology is precisely the topology induced by
$h\mapsto(h,h)$ into $(\cqE,\tau)\times(\cqE,\gamma)$.
Both subspaces in the intersection are closed in this join, so
$C_\cap$ is Hausdorff. The identity on representatives induces
continuous surjections
\begin{equation}
 p_N:C_\cap\longrightarrow\cqQ,\qquad
 p_I:C_\cap\longrightarrow Q_I,\qquad
 \ker p_N=\cqN/H_\cap,\quad
 \ker p_I=I/H_\cap .
 \label{csq:eq-flow-meet-maps}
\end{equation}
The kernel assertions are direct membership tests; continuity follows
from the quotient universal property. These kernel identifications
are also homeomorphisms when $\cqN$ and $I$ have their subspace
topologies from $\tau\vee\gamma$. To prove this, let $q_\cap$ be
the open quotient map and let $A$ denote either subspace.
Because $H_\cap\subset A$,
$q_\cap(A\cap O)=q_\cap(A)\cap q_\cap(O)$ for every open ambient
$O$. Thus its restriction to $A$ is an open quotient map onto
the corresponding kernel with its subspace topology.
The joint map $(p_N,p_I)$ is continuous and injective since its
kernel is zero. No inference of a topological embedding from those
two properties is needed.

Let $\lambda$ be any actual zero of $\Xi$, with its actual
multiplicity $m_\lambda\geq1$. Choose a disk containing no other zero
and retain the full unit
\[
 \Xi(s)=(s-\lambda)^{m_\lambda}A_\lambda(s),
 \qquad A_\lambda(\lambda)\ne0.
\]
For $h\in\cqE$ define
\begin{equation}
 R_\lambda(h)
 =\operatorname{Res}_{s=\lambda}\frac{h(s)}{\Xi(s)}
 =\frac1{2\pi i}\int_{|s-\lambda|=r}
                    \frac{h(s)}{\Xi(s)}\,ds ,
 \label{csq:eq-flow-residue-functional}
\end{equation}
where the circle has positive orientation and lies in this disk.
The integral is independent of the permitted $r$ by Laurent
expansion. It is continuous for $\gamma$: its absolute value is
at most $r\sup_{|s-\lambda|=r}|h(s)|/
\min_{|s-\lambda|=r}|\Xi(s)|$, and a compact supremum is bounded
by a fixed growth seminorm times the maximum of its exponential
weight on the circle. It annihilates $I$, since $h/\Xi$ is entire
there, and therefore annihilates $H_\cap$.

Logarithmic differentiation of the full local factor gives
\begin{equation}
 \frac{c\Xi'(s)}{\Xi(s)}
 =\frac{m_\lambda c}{s-\lambda}
        +c\frac{A_\lambda'(s)}{A_\lambda(s)},\qquad
 R_\lambda(c\Xi')=m_\lambda c .
 \label{csq:eq-flow-logarithmic-unit}
\end{equation}
In particular $\Xi'\notin I$: the displayed pole is nonzero.
Thus $\beta_I=[\Xi']_{Q_I}\ne0$, and
$\widetilde\beta=[\Xi']_{C_\cap}\ne0$ as well, independent of the
status of $\beta$ in the actual source quotient.
The residue functional factors continuously through $C_\cap$; write
its induced functional as $\widetilde R_\lambda$.
The formula
\begin{equation}
 \Pi_\lambda(c)
  =\frac{\widetilde R_\lambda(c)}{m_\lambda}\,
                  \widetilde\beta,\qquad c\in C_\cap
 \label{csq:eq-flow-meet-projection}
\end{equation}
defines a continuous projection onto the line
$\mathbb C\widetilde\beta$. Indeed its square is itself by
$\widetilde R_\lambda(\widetilde\beta)=m_\lambda$, and scalar
multiplication on the quotient is continuous. Its inverse coordinate
on that line is $\widetilde R_\lambda/m_\lambda$, while
$c\mapsto c\widetilde\beta$ is its continuous inverse.
Consequently this line has exactly the ordinary complex topology,
and
$C_\cap=\mathbb C\widetilde\beta\oplus
\ker\widetilde R_\lambda$ is a continuous direct sum.
This proves an information-preserving refinement of the actual
source map, including a continuous recovery coordinate.

The full principal-part map
$h\mapsto(\operatorname{PP}_\lambda(h/\Xi))_{\lambda\in Z(\Xi)}$
has kernel exactly $I$. Vanishing of every negative Laurent
coefficient makes $h/\Xi$ entire at every zero; it is already
holomorphic elsewhere, and the proved entire division bound
puts it in $\cqE$. The reverse inclusion is immediate.
Each finite principal-part coordinate is continuous by the same
circle integral with the appropriate power of $s-\lambda$.
On the velocity line its entire image at every actual zero is
the single term $m_\lambda c/(s-\lambda)$ from
\eqref{csq:eq-flow-logarithmic-unit}, with the unchanged analytic
unit in the regular part.

For general $(\theta,a)$ let
$I_{\theta,a}=V_{\theta,a}I$ and transport
\eqref{csq:eq-flow-meet-quotient} by $V_{\theta,a}$.
Its ambient topology is
$\tau_{\theta,a}\vee\gamma$, since $V_{\theta,a}$ is a
$\gamma$-homeomorphism, and its relation is
$N_{\theta,a}\cap I_{\theta,a}=V_{\theta,a}H_\cap$.
Every map, kernel and projection above is transported exactly.
In particular the transported residue coordinate is
\[
 R_{\lambda;\theta,a}(h)
 =\operatorname{Res}_{s=\lambda}
           \frac{(U_{-\theta}T_{-a}h)(s)}{\Xi(s)}.
\]
For the actual temporal defect it gives
$R_{\lambda;\vartheta,a_j}(\mathcal D_j)=m_\lambda v_j$.
The pair of source and analytic classes therefore has the
explicit injective common lift
$v_j[V_{\vartheta,a_j}\Xi']$ in this transported refinement.
For nonzero heat time this formula retains the inverse heat
operator; it does not replace $V_{\theta,a}I$ by an unproved
ordinary ideal generated by $T_aZ_\theta$.

\subsection{The logarithmic connection and the public blowup residue field}

Take the explicit Newman-time choice $\vartheta(t)=0$.
For the original real flow set
\[
 \Psi_j(t,q,\zeta)=\Xi(\zeta+a_j(t,q)),\qquad
 \mathcal D_j=\partial_t\Psi_j
             =v_j(t,q)\Xi'(\zeta+a_j(t,q)).
\]
At every actual zero $\lambda$ the translated divisor is
$\zeta_{\lambda,j}(t,q)=\lambda-a_j(t,q)$, and the full factorization
is
\[
 \Psi_j=
 (\zeta-\zeta_{\lambda,j})^{m_\lambda}
 A_\lambda(\zeta+a_j).
\]
Thus all multiplicities and the actual analytic unit are retained.
Since $a_j$ is real, its imaginary coordinate is exactly
$\operatorname{Im}\lambda$. Differentiating its location gives
$\partial_t\zeta_{\lambda,j}=-v_j$.

On the complement of the divisor in
$[0,T)\times\mathbb R^3_q\times\mathbb C_\zeta$, define the
complex-valued differential form, smooth in the real variables and
holomorphic in $\zeta$,
\begin{equation}
 \omega_j=\frac{d\Psi_j}{\Psi_j}
 =\frac{\Xi'(\zeta+a_j)}{\Xi(\zeta+a_j)}
                         (d\zeta+da_j),\qquad
 da_j=v_j\,dt+\sum_b\partial_{q_b}a_j\,dq_b .
 \label{csq:eq-flow-logarithmic-connection}
\end{equation}
All label and time coefficients are displayed. Put $s=\zeta+a_j$.
Direct exterior differentiation gives
$d\omega_j=(\Xi'/\Xi)'(s)\,ds\wedge ds+
(\Xi'/\Xi)(s)d^2s=0$.
Equivalently the connection $d-\omega_j$ on the trivial complex
line has zero curvature and annihilates the section $\Psi_j$.
This calculation is on the complement of its zeros and retains
the singular divisor explicitly.

At fixed $(t,q)$ a sufficiently small positively oriented circle
around $\zeta_{\lambda,j}$ has
\[
 \frac1{2\pi i}\int\omega_j=m_\lambda,
\]
where only its $d\zeta$ restriction is integrated. The pole term in
\eqref{csq:eq-flow-logarithmic-unit} gives this integral; the
holomorphic unit term has zero integral by Cauchy's theorem.
The time component and its residue are exactly
\begin{align}
 \omega_j(\partial_t)
 &=\frac{\mathcal D_j}{\Psi_j}
 =\frac{m_\lambda v_j}{\zeta-\zeta_{\lambda,j}}
       +v_j\frac{A_\lambda'(\zeta+a_j)}
                       {A_\lambda(\zeta+a_j)},\notag\\
 r_j(t,q)
 &:=\operatorname{Res}_{\zeta=\zeta_{\lambda,j}}
                    \frac{\mathcal D_j}{\Psi_j}
 =m_\lambda v_j
 =-m_\lambda\partial_t\zeta_{\lambda,j}.
 \label{csq:eq-flow-temporal-residue}
\end{align}
The coordinate change $s=\zeta+a_j$ has differential $ds=d\zeta$
on each spectral fibre, so this is precisely
$R_{\lambda;0,a_j}(\mathcal D_j)$, including its orientation
and multiplicity factor.

Write $r=(r_1,r_2,r_3)^{\mathsf T}$, and retain the full original
polynomial target and its matrix inverse:
\[
 S=(F_1/2,F_2,F_3)^{\mathsf T},\quad
 J=DS=\operatorname{diag}(1/2,1,1)DF,\quad
 J^{-1}=[2W_1\ W_2\ W_3].
\]
The literal recovery and pointwise norm identities are
\begin{equation}
 \begin{aligned}
 u(t,X_t(q))
 &=\frac{2W_1(X_t(q))r_1+
              W_2(X_t(q))r_2+W_3(X_t(q))r_3}{m_\lambda},\\
 |r(t,q)|^2&=m_\lambda^2|J(X_t(q))u(t,X_t(q))|^2 .
 \end{aligned}
 \label{csq:eq-flow-residue-velocity-recovery}
\end{equation}
They follow by $r=m_\lambda v$, $v=Ju$ and the stated matrix
inverse, with no selection of a global inverse sheet of $S$.
The flow $X_t$ is retained as part of the data.

For the nonempty compact velocity support $K$ put
\[
 M_K=\max_K\|J\|_{\mathrm{op}},\qquad
 N_K=\max_K\|J^{-1}\|_{\mathrm{op}}.
\]
They are finite and positive by compactness and $\det J=-1$.
Since $X_t(K)=K$, $X_t$ is onto, and $u,v,r$ vanish outside
the corresponding compact labels, \eqref{csq:eq-flow-residue-velocity-recovery}
gives
\begin{equation}
 \frac{m_\lambda}{N_K}\|u(t)\|_\infty
 \leq\sup_q|r(t,q)|
 \leq m_\lambda M_K\|u(t)\|_\infty .
 \label{csq:eq-flow-residue-supremum}
\end{equation}
The volume-preserving change of variables $x=X_t(q)$ gives
\begin{equation}
 \int_{\mathbb R^3}|r(t,q)|^2\,dq
 =m_\lambda^2\int_{\mathbb R^3}|J(x)u(t,x)|^2\,dx,
 \label{csq:eq-flow-residue-energy}
\end{equation}
whose square root is between
$(m_\lambda/N_K)\|u(t)\|_2$ and
$m_\lambda M_K\|u(t)\|_2$.
These are exact spatial estimates for the arithmetic residues.
With the previously proved spectral normalization
$d_Z(0)^2=\int_{\mathbb R}|\Xi'(s)|^2ds>0$, the unchanged
real translations also give
\[
 |r(t,q)|
 =\frac{m_\lambda}{d_Z(0)}
       \left(\sum_j\|\mathcal D_j(t,q,\cdot)\|_{L^2(\mathbb R)}^2
       \right)^{1/2}.
\]

\begin{cqproposition}[Application of the cited public Navier--Stokes theorem]
For every $\nu>0$, take the velocity, pressure and force supplied by
the public Theorem~1.1 of \cite{openaiNS2026release}. The construction above,
with $\vartheta=0$, gives arithmetic logarithmic residue fields
satisfying
\[
 \sup_{0\leq t<1}\int_{\mathbb R^3}|r(t,q)|^2\,dq<\infty,
 \qquad
 \limsup_{t\uparrow1}\sup_q|r(t,q)|=\infty .
\]
Every residue recovers the full physical velocity by
\eqref{csq:eq-flow-residue-velocity-recovery}, and each translated
arithmetic divisor has the same imaginary coordinates as the
original zero divisor of $\Xi$.
\end{cqproposition}
\begin{proof}
The cited theorem supplies a smooth incompressible velocity on
$[0,1)\times\mathbb R^3$, a fixed compact support, zero initial
velocity, a uniform kinetic-energy bound, and unbounded terminal
velocity supremum, with its stated smooth compactly supported force
and its pressure. These are the inputs of the proved preterminal
flow construction, so its volume and inverse statements apply
on every $[0,T_0]$ with $T_0<1$. Formula
\eqref{csq:eq-flow-residue-energy} and the upper bound using $M_K$
prove the uniform residue energy bound. The lower bound in
\eqref{csq:eq-flow-residue-supremum} proves the unbounded terminal
supremum. Recovery and imaginary-coordinate preservation were
proved above for the actual same fields. The existence step in
this corollary is exactly the cited Theorem~1.1; the correspondence
and every displayed estimate are proved here.
\end{proof}

\subsection{The published witness's exact growing residue}

Take the particular corrected field of the public construction, with
the source's fixed $h\in(0,1/100)$, $X_{\rm in}>0$, $e_0>0$,
and remainder $R_*(\tau)$. Its closing formulas
\cite[equations (10.20)--(10.22), p.~124]{openaiNS2026release}, as imported in \eqref{naw:actualswirl}, are
\[
 |R_*(\tau)|\leq C_*\tau^{2h}\quad(0<\tau<\tau_*),\qquad
 (u_\nu)_2(1-\tau,x_{\nu,\tau})
   =\sqrt\nu\,\tau^{-1/2-h}(e_0+R_*(\tau)),
\]
\[
 r_{\nu,\tau}=\sqrt{2\nu X_{\rm in}\tau},\quad
 x_{\nu,\tau}=(r_{\nu,\tau},0,0),\quad
 q_{\nu,\tau}=X_{\nu,1-\tau}^{-1}(x_{\nu,\tau}).
\]
The remainder is the actual source remainder, including its corrected
fields. The parameter $h$ is the one chosen in that construction.
The quantity $r_{\nu,\tau}$ here is the displayed physical radius;
the map to the original polynomial target is retained below.
The labels $q_{\nu,\tau}$ vary with $\tau$ according to the exact
material inverse just displayed.

Substitution into the full original polynomials gives
\[
 S(r,0,0)=(0,0,2r),\qquad
 J(r,0,0)=
 \begin{pmatrix}
 0&0&1/2\\
 0&1&3r\\
 2&-3r^2&-r^3
 \end{pmatrix}.
\]
Here the same original formulas are
\[
 F_1=Q^3w+y^2Q(4+3xy),\quad
 F_2=y+3xQ^2w+3xy^2(4+3xy),\quad
 F_3=2x-3x^2y-x^3w,\quad Q=1+xy.
\]
Differentiating $S=(F_1/2,F_2,F_3)$ and setting $(x,y,w)=(r,0,0)$
gives each matrix entry. Its componentwise inverse gives
$u_3=2v_1$ and $u_2=v_2-6rv_1$, retaining the off-diagonal $3r$.
The first two arithmetic shifts at $(1-\tau,q_{\nu,\tau})$ are zero.
Consequently, as an identity of full entire functions of $\zeta$,
\begin{equation}
\begin{split}
 \mathcal E_{\nu,\tau}(\zeta)
 &:=\mathcal D_{\nu,2}(1-\tau,q_{\nu,\tau},\zeta)
       -6r_{\nu,\tau}
          \mathcal D_{\nu,1}(1-\tau,q_{\nu,\tau},\zeta)\\
 &=\sqrt\nu\,\tau^{-1/2-h}
                    (e_0+R_*(\tau))\Xi'(\zeta).
\end{split}
 \label{csq:eq-actual-witness-entire-defect}
\end{equation}
Indeed each defect is $v_j\Xi'$ because its shift is zero, so their
combination is $(u_\nu)_2\Xi'$, and the stated source identity applies.

For every actual zero $\lambda$ of multiplicity $m_\lambda$, its
complete local logarithmic image is
\begin{equation}
 \frac{\mathcal E_{\nu,\tau}(\zeta)}{\Xi(\zeta)}
 =\sqrt\nu\,\tau^{-1/2-h}(e_0+R_*(\tau))
     \left(\frac{m_\lambda}{\zeta-\lambda}
                    +\frac{A_\lambda'(\zeta)}{A_\lambda(\zeta)}\right).
 \label{csq:eq-actual-witness-meromorphic-defect}
\end{equation}
Therefore the actual unshifted arithmetic residue at the source
sample is exactly
\begin{equation}
 \operatorname{Res}_{\zeta=\lambda}
       \frac{\mathcal E_{\nu,\tau}(\zeta)}{\Xi(\zeta)}
 =m_\lambda\sqrt\nu\,\tau^{-1/2-h}
                           (e_0+R_*(\tau)).
 \label{csq:eq-actual-witness-residue-growth}
\end{equation}
All multiplicities, the whole analytic unit, the positive viscosity,
the original remainder and the exact material labels remain present.
Choose $\tau_{**}>0$ within the source interval and its cutoff region
with $C_*\tau_{**}^{2h}\leq e_0/2$. For $0<\tau<\tau_{**}$ the
absolute value of \eqref{csq:eq-actual-witness-residue-growth} is at
least $m_\lambda\sqrt\nu\,e_0\tau^{-1/2-h}/2$, which diverges.
This is an explicit pole-coefficient singularity of the published
constructed field's arithmetic image.

For the three residue channels $r_{\nu,j}$ of
\eqref{csq:eq-flow-temporal-residue}, the same evaluation and the
two-component Cauchy--Schwarz inequality give the quantitative bound
\begin{equation}
 |r_\nu(1-\tau,q_{\nu,\tau})|
 \geq
 \frac{m_\lambda\sqrt\nu\,e_0}
             {2\sqrt{1+72\nu X_{\rm in}\tau}}\,
                   \tau^{-1/2-h}.
 \label{csq:eq-actual-witness-residue-vector-growth}
\end{equation}
To verify the denominator, the coefficient vector
$(-6r_{\nu,\tau},1)$ has squared norm
$1+36r_{\nu,\tau}^2=1+72\nu X_{\rm in}\tau$.
Its product with the first two residue coordinates is exactly
the residue in \eqref{csq:eq-actual-witness-residue-growth}; the
third coordinate can only increase the Euclidean norm on the left.
If $E_\nu=\nu^{5/4}F_*(1)$ is the full source force-integral energy
bound and $M_\nu=\max_{K_\nu}\|J\|_{\mathrm{op}}$, the earlier
volume identity simultaneously gives
\[
 \int_{\mathbb R^3}|r_\nu(t,q)|^2dq
 \leq m_\lambda^2M_\nu^2E_\nu^2
 =m_\lambda^2M_\nu^2\nu^{5/2}F_*(1)^2\quad(t<1).
\]

At these particular first two channels, their zero shifts mean the
source quotient is the original $\cqQ$ itself. The exact classes are
\[
 [\mathcal E_{\nu,\tau}]_{\cqQ}
 =\sqrt\nu\,\tau^{-1/2-h}(e_0+R_*(\tau))\beta,\qquad
 [\mathcal E_{\nu,\tau}]_{C_\cap}
 =\sqrt\nu\,\tau^{-1/2-h}(e_0+R_*(\tau))\widetilde\beta.
\]
The latter line has the explicitly proved continuous inverse
$\widetilde R_\lambda/m_\lambda$, so its scalar coordinate has
exactly the displayed singularity. In particular these common
classes leave every bounded subset of $C_\cap$: a continuous linear
functional sends bounded sets to bounded scalar sets, whereas their
$\widetilde R_\lambda$ values diverge. Their projection to the
actual source quotient and its exact kernel remain
\eqref{csq:eq-flow-meet-maps}; the undecided source membership is
not replaced by an identification with the analytic ideal.

The source class has an exact two-case classification as well.
If $\beta=0$, every displayed class in $\cqQ$ is zero.
If $\beta\ne0$, the line $\mathbb C\beta$ has exactly the ordinary
complex topology, and these source classes leave every bounded
subset of $\cqQ$. Here is the complete topological argument,
which does not require a residue functional on $\cqQ$.
The closed relation makes $\cqQ$ a Hausdorff topological vector
space. The compact circle $\{c\beta:|c|=1\}$ omits zero, so it
has a disjoint balanced zero-neighborhood $U$. If $c\beta\in U$
and $|c|\geq1$, multiplying by $1/|c|$ contradicts this
disjointness. Thus $c\beta\in U$ implies $|c|<1$; scaling $U$
proves continuity of the inverse scalar coordinate on the line.
For a bounded subset $B$ there is $t>0$ with $B\subset tU$,
so $c\beta\in B$ implies $|c|<t$. The source lower bound makes
the displayed scalar tend to infinity, proving the asserted escape.
The source membership of $\Xi'$ decides which case occurs; this
argument does not decide that membership. In the same two cases
the full complexified velocity map has respectively rank zero or
rank three, by the exact kernel $J^{-1}(\mathfrak a_\beta^3)$.
There is no intermediate loss of one component in this common-$\beta$
construction.

There is also an exact relation between the changing sample labels
and the time component of the flat logarithmic connection.
Let $A_\tau=D_qX_{\nu,1-\tau}(q_{\nu,\tau})$. Differentiating
$X_{\nu,1-\tau}(q_{\nu,\tau})=x_{\nu,\tau}$ gives
\[
 A_\tau\,\frac{dq_{\nu,\tau}}{d\tau}
 =\frac{dx_{\nu,\tau}}{d\tau}
       +u_\nu(1-\tau,x_{\nu,\tau}),\qquad
 \frac{dx_{\nu,\tau}}{d\tau}
       =\frac{r_{\nu,\tau}}{2\tau}(1,0,0).
\]
The sign follows from $d(1-\tau)/d\tau=-1$.
Since $\nabla_q a_j$ is the $j$th row of $J(x_{\nu,\tau})A_\tau$,
and the first two rows of $J(r,0,0)$ have zero first component,
\[
 \nabla_q a_j\cdot\frac{dq_{\nu,\tau}}{d\tau}=v_j,\qquad j=1,2.
\]
Thus $a_j(1-\tau,q_{\nu,\tau})=0$ has total derivative
$-v_j+\nabla_q a_j\cdot dq_{\nu,\tau}/d\tau=0$, exactly as required.
On the spectral complement of the zeros, the pullback of
\eqref{csq:eq-flow-logarithmic-connection} to this curve and fixed
$\zeta$ is zero in the $\tau$ direction, while its original
$\partial_t$ component has the residue
\eqref{csq:eq-flow-temporal-residue}. These statements are related
by the displayed full tangent-vector identity; the sample is not
silently treated as one fixed material label.

For the full four-term profile $K_\theta=U_\theta K_0$ of
\eqref{nfl:m:K}, the corresponding
actual source classes are instead
\[
 \kappa=[K_0]_{\cqQ},\qquad \beta_K=[K_0']_{\cqQ},\qquad
 [T_aK_\theta]=\overline V_{\theta,a}\kappa,\qquad
 [vT_aK_\theta']=\overline V_{\theta,a}(v\beta_K).
\]
These follow by the same commuting translation, heat and derivative
maps and retain $\kappa$ without assigning it the value zero.
The residue corollary above uses the actual $\Xi$ profile.
The displayed four-term quotient formulas concern the whole $K$
profile and its derivative; they do not identify individual zeros
of $K$ with zeros of $\Xi$.

\subsection{The translated modular double}

Define $C_Eh(s)=\overline{h(\overline s)}$,
$\tau^c=(C_E)_*\tau$ and $N^c=C_E\cqN$.
This is an antilinear involution of the original order class,
since conjugating the Taylor coefficients gives an entire function
with exactly the same maximum modulus on each centered disk.
Also $\|C_Eh\|_{p,a}=\|h\|_{p,a}$, so it is a $\gamma$-homeomorphism.
By definition it is a homeomorphism from $(\cqE,\tau)$ to
$(\cqE,\tau^c)$, with inverse itself. Applying it to the convergent
heat series gives $C_EU_\theta=U_{\overline\theta}C_E$ with the
coefficient $-1/4$ retained. Direct conjugation gives
$C_ET_a=T_{\overline a}C_E$, and the established heat identity gives
$C_EU_\theta=U_{\overline\theta}C_E$.
On the two entire sectors define
\[
 \widehat V_{\theta,a}
   =\operatorname{diag}(T_aU_\theta,T_aU_{-\theta}),\qquad
 \mathcal J(h,k)=(C_Ek,C_Eh).
\]
Then the exact all-parameter identity is
\begin{equation}
 \mathcal J\widehat V_{\theta,a}
 =\widehat V_{-\overline\theta,\overline a}\mathcal J .
 \label{csq:eq-flow-translated-cartan}
\end{equation}
Indeed its first component is
$C_ET_aU_{-\theta}k=T_{\overline a}U_{-\overline\theta}C_Ek$,
and its second component is the same calculation with $\theta$.
For the actual real flow, $a\in\mathbb R$; for real Newman time it
therefore sends $(\theta,a)$ to $(-\theta,a)$.

Give the two ambient sectors the topologies
$(T_aU_\theta)_*\tau$ and $(T_aU_{-\theta})_*\tau^c$, and quotient
by the closed relation
\[
 \widehat N_{\theta,a}
   =T_aU_\theta\cqN\oplus T_aU_{-\theta}N^c .
\]
Applying \eqref{csq:eq-flow-translated-cartan} both to these relation
spaces and to the definitions of their transported topologies proves
an antilinear homeomorphism from this actual doubled quotient to
the one indexed by $(-\overline\theta,\overline a)$.
Its square is the identity, and its conjugacy with the quotient heat
and translation maps follows on representatives from the same
displayed equality. Thus the Cartan reversal retains the full real
flow translation while reversing the radial heat parameter, with
both original source and conjugate-source sectors specified.

\section{The exact strength of the frozen source spectral information}
\label{csq:sec-source-spectral-strength}

The primary spectrum statement in \cite[Proposition~3.6(ii)]{ccmProlate2024}
specifies the actual zero set as the spectrum of the Hadamard
quotient's multiplication operator. It gives no multiplicity
assertion. The cited \cite[Proposition~2.24]{connesMarcolli2008} computes the
finite Hermite polynomial image; the discussion after its proof
explicitly treats function-space choices. The finite multiplicity
calculation in \cite[Theorem~2.47 and equations (2.447)--(2.449)]{connesMarcolli2008}
belongs to the separately specified weighted Hilbert quotient.
The following exact algebra records both the obstruction available
from such data and the limits of transferring a spectrum set alone.

\subsection{Exact algebraic spectral defects of the source relation}
\label{csq:sec-sourcealgebra}
These calculations keep the actual $\cqN$ and require only the
already proved $M\cqN\subseteq\cqN$. For every $\lambda\in\cqC$,
division by $x=s-\lambda$ gives the algebraic identities
\begin{equation}\label{csq:eq-source-kercoker}
 \begin{aligned}
 \operatorname{coker}_{\rm alg}(S-\lambda)
     &=\cqE/(\cqN+x\cqE),\\
 \ker(S-\lambda)&\simeq
        (\cqN\cap x\cqE)/(x\cqN),
        & [n]&\longmapsto[n/x].
 \end{aligned}
\end{equation}
Division belongs to the original order class: if $f(\lambda)=0$,
the filled quotient is entire, is bounded on a fixed disk, and
outside $|s|\geq2|\lambda|+1$ its modulus is at most $2|f(s)|$.
Thus $x\cqE$ is exactly the kernel of evaluation at $\lambda$.
The cokernel identity follows from the image of the induced
multiplication. For the second identity, changing $n$ by $xn_0$
changes $n/x$ by $n_0\in\cqN$; its image is in the stated
kernel. Conversely any kernel representative $f$ has
$xf\in\cqN\cap x\cqE$. Finally $[n/x]=0$ exactly when
$n\in x\cqN$. This proves every asserted kernel and image.

The cokernel has only two possibilities. When every $n\in\cqN$
vanishes at $\lambda$, it is one-dimensional, with its exact
identification given by evaluation. Otherwise retain any
$h_\lambda\in\cqN$ with $h_\lambda(\lambda)\ne0$. Then
\begin{equation}\label{csq:eq-source-division}
 \begin{aligned}
 A_\lambda f(s)&=
  \frac{f(s)-f(\lambda)h_\lambda(s)/h_\lambda(\lambda)}{s-\lambda},\\
 f&=\frac{f(\lambda)}{h_\lambda(\lambda)}h_\lambda
                                      +(s-\lambda)A_\lambda f
 \end{aligned}
\end{equation}
proves surjectivity of $S-\lambda$, and its exact remaining kernel is
\begin{equation}\label{csq:eq-source-splitting}
 \ker(S-\lambda)\simeq
 \cqN/\bigl((s-\lambda)\cqN\oplus\cqC h_\lambda\bigr),
 \qquad n\longmapsto[A_\lambda n].
\end{equation}
The map is onto because a kernel class $[g]$ is the image of
$n=(s-\lambda)g\in\cqN$. Its kernel consists exactly of
$n=(s-\lambda)n_0+c h_\lambda$, $n_0\in\cqN$, by
\eqref{csq:eq-source-division}. The sum is direct by evaluation.
In this second case algebraic invertibility is therefore equivalent
to the exact identity
$\cqN=(s-\lambda)\cqN\oplus\cqC h_\lambda$.
On that locus the representative formula $[f]\mapsto[A_\lambda f]$
is independent of representatives, since $A_\lambda\cqN\subseteq
\cqN$; multiplication in \eqref{csq:eq-source-division} proves the
right inverse, and $A_\lambda((s-\lambda)f)=f$ proves the left
inverse. These are algebraic statements; continuity of this inverse
in $\tau$ has not been inserted.

The first derivative and the heat generator have exactly the same
frozen invariance question for the source subspace:
\begin{equation}\label{csq:eq-source-D-equivalence}
 D\cqN\subseteq\cqN
 \quad\Longleftrightarrow\quad D^2\cqN\subseteq\cqN
 \quad\Longleftrightarrow\quad -\tfrac14D^2\cqN\subseteq\cqN.
\end{equation}
Indeed the product rule is
$D^2(sn)-sD^2n=2Dn$. Since $sn\in\cqN$, second-derivative
invariance puts both terms on the left in $\cqN$, proving
first-derivative invariance. Iteration proves the reverse implication;
the retained nonzero scalar $-1/4$ proves the last equivalence.
If these equivalent inclusions hold, all $D^j\Xi$ belong to
$\cqN$. At each $\lambda$ at least one is nonzero: otherwise
the Taylor series would make $\Xi$ identically zero. Thus every
$S-\lambda$ is then algebraically surjective by
\eqref{csq:eq-source-division}. Its prescribed topological spectral
points would be represented entirely by the kernel
\eqref{csq:eq-source-splitting} or a noncontinuous algebraic inverse.
This computes the precise alternatives without assigning either
one to the source's stated spectral set.

There is also a full multiplicity consequence of algebraic derivative
descent. Writing $d[f]=[Df]$ on its exact descent locus, the product
rule gives $[d,S]=I$. For $0\ne v\in\ker(S-\lambda)$,
\begin{equation}\label{csq:eq-source-weylladder}
 (S-\lambda)d^kv=-k d^{k-1}v,
 \qquad (S-\lambda)^kd^kv=(-1)^k k!v\quad(k\geq0),
\end{equation}
with zero right side in the first identity at $k=0$.
To prove induction, commute the next $d$ using
$(S-\lambda)d=d(S-\lambda)-I$; the coefficient $-k$ becomes
$-(k+1)$. Iteration proves the second identity. A finite linear
dependence among $v,dv,\ldots,d^nv$, with largest nonzero
coefficient at index $j$, becomes that coefficient times
$(-1)^j j!v$ after applying $(S-\lambda)^j$, a contradiction.
Thus every such eigenvector generates an infinite independent
generalized eigenvector sequence. No finite source multiplicity
assertion has been assumed.

Finally even a single forward nonzero heat inclusion has a calculated
consequence. Fix $t\ne0$ and $c=t/2$, and define
\begin{equation}\label{csq:eq-source-heatcommutator}
 \begin{gathered}
 A_0=U_t,\qquad A_{k+1}=MA_k-A_kM,\qquad A_k=P_k(D)U_t,\\
 P_0(X)=1,\qquad P_{k+1}(X)=cX P_k(X)-P_k'(X).
 \end{gathered}
\end{equation}
The identity follows by induction from $[M,U_t]=cDU_t$ and
$[M,D^j]=-jD^{j-1}$: commuting $M$ through $P_k(D)U_t$
contributes $-P_k'(D)U_t+cP_k(D)DU_t$.
The degree of $P_k$ is exactly $k$, with leading coefficient
$c^k\ne0$. Under $U_t\cqN\subseteq\cqN$, induction with
$M\cqN\subseteq\cqN$ gives $A_k\cqN\subseteq\cqN$;
subtracting the lower-degree terms of $P_k$ then gives exactly
\begin{equation}\label{csq:eq-source-heatconsequence}
 D^kU_t\cqN\subseteq\cqN\quad(k\geq0).
\end{equation}
In particular every derivative of the nonzero $Z_t=U_t\Xi$ lies
in $\cqN$. At every point one derivative is nonzero by the entire
Taylor theorem, so every $S-\lambda$ is algebraically surjective
in this case too. This proof uses only the one forward inclusion.
For the stronger equality $U_t\cqN=\cqN$, it also proves
$D\cqN\subseteq\cqN$, by writing each $n\in\cqN$ as
$U_tm$ in \eqref{csq:eq-source-heatconsequence}. None of these
exact implications presumes the inclusion or equality for the actual
source closure. The explicit obstruction spaces above remain the
calculated source question.

\subsection{One-sided heat inclusion and the full generalized chain}

The already proved commutator identity
\eqref{csq:eq-source-D-equivalence} and its following heat calculation
show that $D\cqN\subseteq\cqN$ makes every $S-\lambda$ algebraically
onto. The single inclusion $U_t\cqN\subseteq\cqN$ at any $t\ne0$
also makes every $S-\lambda$ onto, using every derivative
of $U_t\Xi$ and the exact source division map.
Here is a stronger consequence that needs no descended derivative.

\begin{cqlemma}
Let $T:V\to V$ be a surjective endomorphism of a complex vector
space and let $0\ne v_0\in\ker T$. Then for every integer $k\geq1$,
$\dim\ker T^k\geq k$. More precisely one can choose
$Tv_j=v_{j-1}$ for every $j\geq1$, and the vectors
$v_0,\ldots,v_n$ are linearly independent for every $n$.
\end{cqlemma}
\begin{proof}
Surjectivity gives the successive preimages. It follows inductively
that $T^jv_j=v_0$ and $T^{j+1}v_j=0$. Applying $T^n$ to a relation
$\sum_{j=0}^n c_jv_j=0$ gives $c_nv_0=0$, hence $c_n=0$.
Repeating with the remaining highest index gives every coefficient
zero. The first $k$ vectors lie in $\ker T^k$, proving the bound.
\end{proof}

Applied to $T=S-\lambda$, this proves that either frozen
inclusion just stated is incompatible with any nonzero
finite-dimensional generalized eigenspace of the actual source
operator. The source-independent implication is
\begin{equation}
\begin{gathered}
 U_t\cqN\subseteq\cqN,\quad t\ne0,\quad
                  0\ne v_0\in\ker(S-\lambda)
 \quad\Longrightarrow\quad
 \dim\ker(S-\lambda)^k\geq k\quad(k\geq1).
\end{gathered}
\label{csq:eq-source-surjective-chain}
\end{equation}
With a descended $d[f]=[Df]$, its more specific ladder is
$v_j=(-1)^jd^jv_0/j!$; indeed
$[d,S]=I$ gives $(S-\lambda)d^jv_0=-j d^{j-1}v_0$.
Thus the complete factorial and sign agree with the general
preimage construction. This is a proved algebraic incompatibility;
it does not assert that the unnamed source quotient has the
finite generalized block needed to apply the incompatibility.

\subsection{The source Sobolev cokernel and an exact common test-space map}
\label{csq:sec-sobolev-comparison}
The Sobolev construction cited immediately after Connes--Marcolli
Proposition 2.24 is specified in Chapter 2, Section 9. It is a different,
explicitly defined quotient, and supplies a useful exact comparison.
We retain its weight, its generator, and its limited set of detected jets.
The primary locators are Connes--Marcolli, printed pp.407--408 and
416--420 \cite{connesMarcolli2008}, and Connes, Section III, pp.10--15 and
Appendix I, pp.62--67 of the author PDF \cite{connesSelecta1999}.
Here and below $\delta>1$.

For the rational field, choose the splitting
$C_{\mathbb Q}=C_{\mathbb Q,1}\times\mathbb R_+^*$ used in the source,
give the compact factor Haar mass one, and put $x=e^u$ on the second
factor. Its trivial-character component is the Hilbert space
\[
 B_\delta=L^2\bigl(\mathbb R,(1+u^2)^{\delta/2}\,du\bigr),
 \qquad
 \|f\|_\delta^2=\int_{\mathbb R}|f(u)|^2(1+u^2)^{\delta/2}\,du.
\]
The source adelic map is
\[
 E_{\rm ad}\psi(g)=|g|^{1/2}\sum_{q\in\mathbb Q^*}\psi(qg),
 \quad
 \psi\in\mathcal S(\mathbb A_{\mathbb Q})_0
 =\{\psi:\psi(0)=0,\ \int\psi\,dx=0\}.
\]
Let $P_1$ denote averaging over $C_{\mathbb Q,1}$, and define
\[
 R_\delta=\overline{P_1E_{\rm ad}
             (\mathcal S(\mathbb A_{\mathbb Q})_0)}^{\,B_\delta},
 \qquad H_{\delta,1}=B_\delta/R_\delta.
\]
These are exactly the trivial-character relation space and cokernel
of the source. Indeed the source completion uses the norm of its image,
so the extended isometry has closed range equal to the closure of the
uncompleted image. Compact averaging commutes with that equivariant
closed range. If $b$ is a limit of image vectors and satisfies $P_1b=b$,
averaging an approximating sequence gives a sequence in the averaged
image converging to $b$. This proves that the displayed closure is the
trivial-character part of the full relation space, without replacing
any source relation by an orthogonal complement as a definition.

We import explicitly the following result of the source proof:
under the bilinear pairing $\langle f,\eta\rangle=\int f(u)\eta(u)\,du$,
the annihilator of $R_\delta$ in
$B_{-\delta}=L^2(\mathbb R,(1+u^2)^{-\delta/2}\,du)$ is the closed
linear span of
\begin{equation}\label{csq:eq-sobolev-source-annihilator}
 u^a e^{itu},\qquad
 \zeta(\tfrac12+it)=0,\quad t\in\mathbb R,\quad
 0\leq a<m_{-t},\quad a<\frac{\delta-1}{2}.
\end{equation}
Here $m_{-t}$ is the original multiplicity of the zero $-t$ of
$\Xi$, equal to that of $\tfrac12+it$ by the retained reflected
Mellin coordinate. The source proves the finite combinations after a
compact Fourier cutoff and proves that the cutoff vectors are dense
by uniformly bounded approximate units. Thus the closed span, not
only its finite-dimensional projections, is the imported statement.
No such annihilator statement is imported for the unnamed topology
$\tau$ of the entire-function ring.

The precise cutoff can also be checked directly from the defined norm:
\[
 \|u^a e^{itu}\|_{-\delta}^2
   =\int_{\mathbb R}|u|^{2a}(1+u^2)^{-\delta/2}\,du<\infty
 \quad\Longleftrightarrow\quad 2a-\delta<-1.
\]
Near zero the integrand is locally integrable for every integer $a\geq0$;
for $|u|\geq1$ it is bounded above and below by positive constants
times $|u|^{2a-\delta}$. Integration of this power proves the strict
inequality and the logarithmic divergence at equality. Hence set
\begin{equation}\label{csq:eq-sobolev-jetcount}
 q_\delta=\left\lceil\frac{\delta-1}{2}\right\rceil,
 \qquad n_{\delta,\lambda}=\min(m_\lambda,q_\delta).
\end{equation}
There are exactly $n_{\delta,\lambda}$ allowed degrees at a real zero
$\lambda$. The original Connes Theorem 1, Section III, p.13, states
the equivalent multiplicity bound $n<(1+\delta)/2$, $n\leq m_\lambda$.
The inspected Connes--Marcolli printed p.408 instead displays
$n<1+\delta/2$ in Theorem 2.47; its proof, pp.419--420,
equations (2.444) and (2.448), gives precisely the stricter bound
proved above. This records the primary discrepancy rather than
silently identifying the two printed formulas. No assumption that
an actual zeta zero has multiplicity greater than one is used.

Let $\mathcal E=C_c^\infty(\mathbb R;\mathbb C)$ retain its usual
fixed-support smooth-function inductive-limit topology, and let
$\mathscr Ff(s)=\int f(u)e^{-ius}\,du$ as above. Define
\begin{align}
 K_{\rm full}
   &=\{f\in\mathcal E:\ D^a\mathscr Ff(\lambda)=0
               \text{ for every zero }\lambda\text{ of }\Xi,
                         \ 0\leq a<m_\lambda\},
       \label{csq:eq-sobolev-full-kernel}\\
 K_\delta
   &=\{f\in\mathcal E:\ D^a\mathscr Ff(\lambda)=0
               \text{ for every real zero }\lambda\text{ of }\Xi,
                         \ 0\leq a<n_{\delta,\lambda}\}.
       \label{csq:eq-sobolev-kernel}
\end{align}
The already proved entire division theorem identifies
$K_{\rm full}=\{f:\mathscr Ff\in\Xi\cqE\}$ exactly.

\begin{cqtheorem}[Exact dense test ingress into the source Sobolev cokernel]
\label{csq:thm-sobolev-ingress}
The map $j_\delta:\mathcal E\to H_{\delta,1}$ sending a test function
to its Hilbert quotient class is continuous, has dense image, and
has kernel exactly $K_\delta$. It induces the continuous map
\begin{equation}\label{csq:eq-sobolev-quotient-map}
 J_\delta:\mathcal E/K_{\rm full}\longrightarrow H_{\delta,1},
 \qquad [f]\longmapsto[f],
 \qquad \ker J_\delta=K_\delta/K_{\rm full}.
\end{equation}
Its image is dense. These statements neither assert surjectivity
nor identify the topology of its image with a quotient topology.
For the source translation generator $D_1$, this ingress retains
the original spectral coordinate through
\begin{equation}\label{csq:eq-sobolev-generator}
 j_\delta(Tf)=iD_1j_\delta(f),
 \qquad T=-i\partial_u,\qquad \mathscr F(Tf)=s\mathscr Ff.
\end{equation}
\end{cqtheorem}
\begin{proof}
The inclusion $\mathcal E\to B_\delta$ is continuous on every
fixed-support space: on $[-L,L]$,
\[
 \|f\|_\delta^2\leq
   \left(\int_{-L}^L(1+u^2)^{\delta/2}\,du\right)\sup|f|^2.
\]
The inductive-limit universal property and the continuous Hilbert
quotient map give continuity of $j_\delta$.
Smooth compactly supported functions are dense in $B_\delta$:
truncate an arbitrary $f$ outside $[-L,L]$, whose squared norm tail
tends to zero by integrability; on the enlarged interval
$[-L-1,L+1]$, the weight is bounded above and below by positive
constants, and convolution with a smooth compactly supported
approximate identity approximates the truncation in ordinary $L^2$
and therefore in the weighted norm. These convolutions have compact
support and are smooth. Consequently their quotient images are dense.

The continuous dual of $B_\delta$ is exactly $B_{-\delta}$ under
the displayed bilinear pairing. To check this, multiply by
$(1+u^2)^{\delta/4}$ to identify $B_\delta$ isometrically with
ordinary $L^2$; the Riesz representation there, with complex
conjugation absorbed in the representing function, gives exactly
the inverse-weight dual and its norm. Since $R_\delta$ is closed,
a vector belongs to it exactly when every member of its annihilator
vanishes on that vector. One direction is the definition. For the
other, a vector outside a closed Hilbert subspace has a nonzero
orthogonal component, and its inner product defines a continuous
separating functional. By the explicitly imported closed-span
description \eqref{csq:eq-sobolev-source-annihilator}, it is enough
to test the listed log-monomial characters. But
\[
 D^a\mathscr Ff(-t)=(-i)^a\int u^a e^{itu}f(u)\,du,
 \qquad
 \int u^a e^{itu}f(u)\,du=i^aD^a\mathscr Ff(-t).
\]
Every factor $i^a$ is nonzero. The coordinate $\lambda=-t$ and
the proved count \eqref{csq:eq-sobolev-jetcount} now give
$\ker j_\delta=K_\delta$. Every equation defining $K_\delta$
appears among the equations defining $K_{\rm full}$, proving
$K_{\rm full}\subseteq K_\delta$. Thus
\eqref{csq:eq-sobolev-quotient-map} is well defined and continuous
by the quotient universal property; changing the representative
changes its image by zero exactly for a member of $K_\delta$.
This proves the kernel and the dense-image assertion.

Finally the source left translation is
$(V(e^\varepsilon)f)(u)=f(u-\varepsilon)$, so its infinitesimal
generator on $\mathcal E$ is $-\partial_u$. The difference quotient
converges in $B_\delta$ because the supports remain in a fixed
compact set and the smooth-function difference quotient converges
uniformly there. Passing to the quotient proves that
$j_\delta(f)$ belongs to the domain of $D_1$ and
$D_1j_\delta(f)=j_\delta(-f')$. Multiplication by $i$ gives the
first equation of \eqref{csq:eq-sobolev-generator}. Integration by
parts in the compactly supported Fourier integral gives the second.
In particular the source eigenparameter $it$ becomes $-t$ under
$iD_1$, exactly the same $s$ coordinate used in the jet equations;
no sign change is absorbed into the source operator.
\end{proof}

For completeness the unspecified source topology also has an exact
comparison, requiring no unproved inclusion of its kernel in
$K_\delta$. Retain
\[
 K_\tau=\{f\in\mathcal E:\mathscr Ff\in\cqN\},
 \qquad L_{\tau,\delta}=K_\tau\cap K_\delta.
\]
Give $\mathcal E$ the join of its usual test topology, the
inverse-image topology $\mathscr F^{-1}\tau$, and the seminorm
$\|\cdot\|_\delta$. The last seminorm is already continuous in
the test topology, but is displayed to specify every target map.
Then $K_\tau$, $K_\delta$ and $L_{\tau,\delta}$ are closed, and
there are continuous maps
\begin{equation}\label{csq:eq-sobolev-source-span}
 \cqQ\ \longleftarrow\
 \mathcal E/L_{\tau,\delta}
 \ \longrightarrow\ H_{\delta,1},
 \qquad
 [f]\longmapsto[\mathscr Ff],\quad [f]\longmapsto[f].
\end{equation}
Their respective kernels are exactly
$K_\tau/L_{\tau,\delta}$ and $K_\delta/L_{\tau,\delta}$;
the right image is dense. Indeed each composite on $\mathcal E$
is continuous by definition, vanishes on the intersection, and
has the asserted kernel before quotienting. The quotient universal
property proves the assertions, and density follows from the same
set of compact tests as above. The map to the product of the two
targets is injective because its kernel is precisely the
intersection already divided out. No topological embedding of that
image is inferred from injectivity. This is a proved relation to
the actual source quotient while its named topology remains retained.
The comparison also intertwines the coordinate operators: $T$
preserves $K_\tau$ because $\mathscr F(Tf)=s\mathscr Ff$ and
$M\cqN\subset\cqN$; it preserves $K_\delta$ because
\[
 D^a(sF)(\lambda)=\lambda D^aF(\lambda)
                              +aD^{a-1}F(\lambda),
\]
where the second term is zero for $a=0$ and otherwise uses an
already imposed lower jet. Thus $T$ descends to the common span,
whose left map intertwines it with $S$ and whose right map
intertwines it with $iD_1$ on the proved domain of the ingress.

There is also a literal compatibility of the original finite Hermite
traces with the adelic map. For the source's even $\psi$, put
$\psi_{\rm ad}=\psi\otimes\mathbf1_{\widehat{\mathbb Z}}$.
It lies in $\mathcal S(\mathbb A_{\mathbb Q})_0$ because its value
and additive integral are the two original zero constraints.
On the idele represented by $(x,1,1,\ldots)$, a rational number
belongs to every $\mathbb Z_p$ exactly when it is an integer:
write it in lowest terms and use any prime dividing its denominator
for the reverse implication. Therefore
\[
 E_{\rm ad}\psi_{\rm ad}(x)
   =x^{1/2}\sum_{n\in\mathbb Z\setminus\{0\}}\psi(nx)
   =2\mathcal E\psi(x).
\]
Every norm-one rational idele class has a representative consisting
of a real sign and finite units. Indeed, for an idele $g$ put
$q=\prod_p p^{v_p(g_p)}\in\mathbb Q_{>0}$, a finite product.
Then $g_p/q\in\mathbb Z_p^*$ for every prime, while the
norm-one identity gives $|g_\infty|=q$. Dividing by the diagonal
rational $q$ gives the asserted representative. The function
is unchanged by the finite units; evenness removes the real sign.
It therefore lies in the entire trivial-character component.
The Fourier transform in $u=\log x$ is thus exactly twice the
original source trace, retaining the factor $2$. In particular the
finite Hermite generators yield $2R_\ell^\pm(s)\Xi(s)$ in this
explicit realization. Their logarithmic representatives lie in
every $B_\delta$ by their proved rapid decay. This proves a concrete
arithmetic compatibility as well as the quotient comparison; it
does not identify a Sobolev topology with the topology on all of
$\cqE$.


\subsection{The actual finite block in the specified Sobolev source}

Retain the source's real frequency $t$, compact character $\chi_0$,
weight $\delta>1$, and actual zero multiplicity $m$.
The vectors in the dual block of the specified source are
\[
 \eta_{t,a}(x)=\chi_0(x)|x|^{it}(\log|x|)^a,\qquad
 a\in\mathbb Z_{\geq0},\quad a<m,\quad a<(\delta-1)/2 .
\]
The last strict bound is exactly integrability of
$u^{2a}(1+u^2)^{-\delta/2}$ on $\mathbb R$: at infinity
its power is $2a-\delta$, integrable precisely when it is
less than $-1$; near zero it is locally integrable for all
these $a$. Let $r$ be the number of permitted integers.
The source representation at $\lambda>0$ acts by
\[
 \vartheta_m^{\rm tr}(\lambda)\eta_{t,a}
 =\lambda^{it}\sum_{b=0}^a\binom ab(\log\lambda)^b\eta_{t,a-b}.
\]
Differentiating at $\lambda=e^\epsilon$, $\epsilon=0$, proves
\begin{equation}
 (D_{\chi_0}^{\rm tr}-it)\eta_{t,a}=a\eta_{t,a-1},\qquad
 (D_{\chi_0}^{\rm tr}-it)^k\eta_{t,a}
 =\frac{a!}{(a-k)!}\eta_{t,a-k}\quad(0\leq k\leq a).
 \label{csq:eq-source-Sobolev-Jordan-block}
\end{equation}
For $k>a$ the expression is zero. These vectors are independent:
divide a proposed relation by the nonzero character and use
that $\log|x|$ ranges over $\mathbb R$, making the remaining
polynomial identically zero. Thus the block has dimension
$r$, one-dimensional ordinary eigenspace, and nilpotence
index $r$. Its multiplicity is generalized multiplicity.
The common-test morphism to the actual Hadamard source is
\eqref{csq:eq-sobolev-source-span}, with its exact two kernels; this finite
block is not silently identified with a block in that
different quotient.

\subsection{An explicit spectral model and the inverse at zero}

Let $\Lambda=Z(\Xi)$ be the actual zero set. On the complex
vector space
\[
 V_\Lambda=\bigoplus_{\lambda\in\Lambda}
                   \bigoplus_{n\geq0}\mathbb C e_{\lambda,n}
\]
every vector has finite support. Give it the finest locally
convex topology, defined by all seminorms on this vector space.
Coordinate functionals separate points, so it is Hausdorff.
Every linear map from it to any locally convex space is
continuous: each pullback of a continuous seminorm is among
these defining seminorms.
With $e_{\lambda,-1}=0$ set
\[
 S_\Lambda e_{\lambda,n}=\lambda e_{\lambda,n}+e_{\lambda,n-1},
 \qquad d_\Lambda e_{\lambda,n}=-(n+1)e_{\lambda,n+1}.
\]
Both are continuous, everywhere-defined, and direct evaluation
on the basis proves $[d_\Lambda,S_\Lambda]=I$.
For $\mu\notin\Lambda$ the full continuous inverse is
\begin{equation}
 (S_\Lambda-\mu)^{-1}e_{\lambda,n}
 =\sum_{j=0}^n
       \frac{(-1)^j}{(\lambda-\mu)^{j+1}}e_{\lambda,n-j}.
 \label{csq:eq-source-model-resolvent}
\end{equation}
Multiplication in either order cancels adjacent terms of
this finite sum and leaves the basis vector. The formula
preserves finite support, so it defines an inverse on
the whole space, continuous by its topology.
For $\mu\in\Lambda$, the $\mu$-summand is a surjective
backward shift and every other summand uses the same finite
inverse formula. Thus $S_\Lambda-\mu$ is onto, its kernel
is exactly $\mathbb C e_{\mu,0}$, and
\begin{equation}
 \ker(S_\Lambda-\lambda)^k
   =\operatorname{span}\{e_{\lambda,0},\ldots,e_{\lambda,k-1}\},
 \quad
 \operatorname{Spec}_{\rm alg}S_\Lambda
  =\operatorname{Spec}_{\rm end}S_\Lambda
  =\operatorname{Spec}_{\rm point}S_\Lambda=\Lambda .
 \label{csq:eq-source-model-spectrum}
\end{equation}
Here $\operatorname{Spec}_{\rm end}$ means failure of a
continuous two-sided inverse in the algebra of continuous
endomorphisms, not the Hilbert spectral subdivision called
continuous spectrum. Since $\Xi(0)\ne0$, this $S_\Lambda$
has a continuous inverse at zero.

There is an exact principal-part realization of the model:
send $e_{\lambda,n}$ to the class of $w^{-n-1}$ in
$\mathbb C[w,w^{-1}]/\mathbb C[w]$ in its $\lambda$-summand.
The finite negative Laurent coefficients give its inverse.
Multiplication by $\lambda+w$ is $S_\Lambda$ and
differentiation in $w$ is $d_\Lambda$, as evaluation on
each monomial proves.
This complete model shows that the correct zero set as
spectrum, geometric multiplicity one, and a continuous
inverse at zero can coexist with a continuous Weyl pair.
It does not identify this model with $\cqQ$.

For the actual quotient, injectivity of $S$ gives exactly
$\cqN\cap s\cqE=s\cqN$: if $sf\in\cqN$, then
$S[f]=0$, so $f\in\cqN$; the converse is polynomial invariance.
The original source inverse at zero is explicitly
\begin{equation}
 S^{-1}[f]=\left[
 \frac{f(s)-f(0)\Xi(s)/\Xi(0)}{s}\right].
 \label{csq:eq-source-explicit-zero-resolvent}
\end{equation}
The singularity is removable, and division by a degree-one
polynomial preserves the original order class. Multiplication
by $s$ gives $[f]$, and injectivity gives its unique inverse.
If $f$ is changed by $n\in\cqN$, the numerator difference
belongs to $\cqN\cap s\cqE=s\cqN$, proving representative
independence directly. The usual topological reading of
the source resolvent also supplies continuity of this
quotient inverse, not continuity of its chosen lift.
More generally,
\[
 S-\lambda\text{ injective}\quad\Longleftrightarrow\quad
 \cqN\cap(s-\lambda)\cqE=(s-\lambda)\cqN
\]
by the same calculation. Saturation at zero does not prove
saturation at other points.
The missing frozen-descent decision can therefore be
approached through an actual finite generalized source block
or a range obstruction, as well as through identification
of the topology. Neither has been supplied by the primary
Hadamard spectrum clauses inspected here.

